\scanpage{10}
Soit $G'$ défini par \deleted{\unsure{le}}la condition d'avoir un diagramme cartésien
\diagram{\begin{tikzcd}[ampersand replacement=\&,column sep=large,execute at end picture={\draw (middle.north west) rectangle (rightnode.south east);}]
\unsure{\mathscr X|X|\mathfrak X}=G\arrow[d,"f"'] \& |[alias=middle]|\mathscr X'=G'\arrow[l]\arrow[d,"f'"]\arrow[dl,phantom,"\text{\uline{cart}}" description] \& |[alias=rightnode]|\mathscr B=\widetilde{\mathscr X}\arrow[l,"p"']\arrow[dl,"{\unsure{\widetilde f'|\widetilde f}}"]\\
I \& I'=T\arrow[l,"\pi"'] \&
\end{tikzcd}}
Utilisant le fait que $f$ est normal, donc $f'$ normal, on trouve que \deleted{\unsure{$f$}} $G'$ est normal.

On va définir un hom
\[
p\colon\mathscr B\longrightarrow G'
\]
i.e. des hom
\[
\widetilde f\colon\mathscr B\longrightarrow T,\quad
\deleted{\unsure{u}}\underline h\colon\mathscr B\longrightarrow G
\qquad\text{tels que }\pi\widetilde f=fh
\]

On prendra
\[
\left\{\begin{aligned}
\underline h(g,B)&=g\qquad \text{i.e. }\underline h\text{ est le composé }\mathscr B\subset G_D=G\times_S D\xrightarrow{\mathrm{pr}_1}G\\
\underline{\widetilde f}(g,B)&=\dot g\in B/B_u(-)\simeq T(-),\quad\text{i.e. }\widetilde f\text{ est le composé}\\
&\hspace{12mm}\mathscr B\longrightarrow\mathscr B/\mathscr B_u=T_D\simeq T\times_S D\xrightarrow{\mathrm{pr}_1}T
\end{aligned}\right.
\]

Pour voir que $\pi\widetilde f=fh$, il suffit de le voir sur l'ouvert dense \inserted{$\mathscr B^{rs}$} de $\mathscr B$ des \deleted{\unsure{po|pa}} $(g,B)$ avec $g$ régulier ss, \deleted{\unsure{où|si on}} \unsure{donc|d'où} $g\in T\subset B$, et alors c'est trivial.

\uline{NB}\quad $\left\{\begin{aligned}
&h\text{ est propre, donc }p\text{ est propre (on va voir qu'il est birationnel)}\\
&\widetilde f\text{ est lisse.}
\end{aligned}\right.$

$G^{\mathrm{reg}}$ l'ouvert des points réguliers de $G=$ ouvert des pts de lissité de $f$

$G'^{\mathrm{reg}}$ son image inverse dans $G'$

$\mathscr B^{\mathrm{reg}}$ son image inverse dans $\mathscr B$.
% archive sheet number: 9
